\unnumberedchapter{Заключение}

В~работе оптимизирован архиватор MZ77, основанный на алгоритме LZ77.
Каждый шаг оптимизации измерялся и профилировался на двух компьютерах с
процессорами AMD Zen~2 и Intel Tiger Lake, а формат архива не менялся:
архивы оптимизированной программы совпадают байт в байт с архивами
исходной.

Основные результаты работы состоят в следующем:
\begin{enumerate}
\item Профилирование исходной программы показало три источника потерь:
сдвиг упреждающего буфера функцией \texttt{memmove} после каждого токена,
побайтовое сравнение кандидатов через функцию \texttt{byte\_at} и вызов
\texttt{fputc}/\texttt{fgetc} на каждый байт архива и результата.
\item Оценена сложность исходной программы: сжатие
$O\bigl(n \cdot (C \cdot M + L)\bigr)$, распаковка $O(n)$, память $O(1)$.
На реальных данных преобладает член $O(T \cdot L)$ "--- около 2200 байт
копирования на токен.
\item Алгоритмическая оптимизация "--- общий буфер окна и упреждающего
буфера "--- исключила этот член и ускорила сжатие в 1,29--1,52~раза.
\item Четыре машинно-независимые оптимизации: блочный ввод-вывод
(распаковка в 1,63--1,67~раза быстрее), встраивание побитового
ввода-вывода (ещё в 1,29--1,43~раза), сравнение по 8~байт (сжатие в
1,66--1,76~раза быстрее "--- самый большой выигрыш работы), оптимизация при
компоновке и по профилю (ещё 5--7~\%).
\item Машинно-зависимые оптимизации: сравнение командами AVX-512BW
ускорило сжатие на процессоре Intel в 1,21~раза, командами AVX2 на
процессоре AMD "--- в 1,13~раза. Реализация выбирается во время выполнения,
поэтому программа работает на любом процессоре x86-64.
\item Оценена сложность оптимизированной программы: сжатие
$O\bigl(n \cdot C \cdot \lceil M/w \rceil\bigr)$, распаковка $O(n)$,
память $O(1)$. Асимптотика осталась линейной, а константа уменьшилась;
измерение подтверждает, что число команд на байт не зависит от размера
входа и уменьшилось в 1,7--4,2~раза.
\end{enumerate}

В~итоге сжатие стало быстрее в 3,1~раза на компьютере с процессором AMD
(с 10,1 до 3,3~с на корпус 291~МБ) и в 2,9~раза на компьютере с
процессором Intel, распаковка "--- в 2,4~раза на обоих.

Работа показала, что наибольший выигрыш дают изменения самого кода,
найденные по профилю: сдвиг буфера и побайтовое сравнение вместе ускорили
сжатие в 2,3--2,5~раза, тогда как флаги компилятора, LTO и PGO "--- лишь на
несколько процентов. Машинно-зависимые оптимизации нужно проверять на том
процессоре, для которого они написаны: на Tiger Lake более узкое
сравнение AVX2 оказалось на 3~\% быстрее, чем AVX-512, а на Zen~2 команды
BMI2 \texttt{pext}/\texttt{pdep} пришлось отвергнуть из-за их стоимости.

Возможные продолжения работы: выбирать AVX2 вместо AVX-512 на процессорах
Intel, где он быстрее; хранить в цепочках хеш-таблицы 32-битные позиции
вместо 64-битных, чтобы уменьшить объём таблиц; исследовать причину
замедления реализации AVX2 на файлах с короткими совпадениями.
